\chapter{Macros}
\label{ch:macros}

Chapter~\ref{ch:language} introduced the period command as a way of repeating the single most recent change. A macro generalizes that idea from one change to an arbitrary recorded sequence of them, and Part VI opens with it because a macro, more than anything else in this book, sits at the boundary between editing and programming: it is composed entirely of commands already covered in Parts II through V, yet using one well is closer to writing a small, disposable program than to issuing a single command.

\section{Recording}

\texttt{q}\textit{letter}, pressed in Normal mode, begins recording into register \textit{letter} (the same register namespace covered in Chapter~\ref{ch:registers}); every keystroke from that point forward is recorded, verbatim, until \texttt{q} is pressed again to stop. \texttt{@}\textit{letter} then replays the recorded sequence, and \texttt{@@} repeats whichever macro was most recently played, without needing to specify its register name again. A count prefixed to \texttt{@} replays the macro that many times in succession: \texttt{10@q} runs the macro in register \texttt{q} ten times.

What is recorded is not a semantic description of an edit but the literal keystroke sequence that produced it: an operator, a motion or text object, perhaps a search, perhaps an Insert-mode segment terminated by escape, exactly as it was typed. This is worth stating plainly because it explains both a macro's chief strength and its chief limitation. The strength is that a macro can capture anything expressible in Vim's own grammar, with no separate scripting syntax to learn beyond what Parts II through V have already covered. The limitation is that a macro replays the same keystrokes against whatever the cursor and buffer happen to look like at replay time, not the same literal effect: a macro recorded as "delete the word under the cursor" will delete whatever word the cursor happens to be on when replayed, which may or may not be the word the macro's author had in mind when validating it against the first instance.

\section{Testing Before Scaling}

The workflow that follows from this limitation, and that experienced Vim users converge on independently, is to test a macro cautiously before trusting it broadly. Record the transformation once, carefully, against a single representative instance; replay it once or twice more with \texttt{@q} to confirm it behaves as intended against a few more instances; and only once it has proven reliable, replay it at scale with a count, \texttt{100@q} or more, across the remainder of the buffer. This is not merely caution for its own sake. A macro that fails partway through a large replay count typically does so because some instance among the targets differs structurally from the one the macro was recorded against, and catching that failure after two or three replays, while it is still easy to identify which instance broke the pattern, is considerably less costly than catching it after the macro has already run one hundred times and left the buffer in a partially, inconsistently transformed state.

This staged approach also gives a concrete, practical shape to a more general question worth asking before recording any macro at all: is this transformation going to be repeated often enough, across similar enough material, to be worth the overhead of recording and validating it in the first place. A transformation needed three times is very often faster to perform by hand three times than to record, test, and replay; a transformation needed fifty times, across material that is genuinely structurally consistent, is a clear case for a macro. The deciding factor is not the raw repetition count alone but repetition count weighed against how reliably the surrounding structure repeats: a macro applied against material with inconsistent formatting can cost more in cleanup after a failed run than it would have cost to simply perform the edits by hand from the start.

\section{Nested and Recursive Macros}

A macro recorded into one register can invoke another by including \texttt{@}\textit{letter} as part of its own recorded keystrokes, which lets simple, single-purpose macros compose into more elaborate ones exactly as operators compose with motions in Part II: rather than recording one long, monolithic macro, several small macros, each handling one well-defined transformation, can be built and validated independently, then combined by having one invoke another.

A macro that invokes itself, by including \texttt{@q} as part of what is recorded into register \texttt{q}, becomes recursive, repeating until it encounters a condition under which the step it performs (frequently a search) fails, at which point the recursive call itself fails and the chain of replays stops. This is a genuine, if minimal, form of conditional looping, expressed entirely within Vim's own Normal-mode grammar rather than requiring the Vimscript control-flow constructs covered in Chapter~\ref{ch:vimscript}: a recursive macro that searches for a pattern, transforms the found instance, and then invokes itself again will continue until no further instance of the pattern remains to be found, at which point the search fails, the recursive invocation fails along with it, and the macro terminates on its own.

\section{Defining a Macro Without Recording It}

A macro's contents are, ultimately, just the text stored in a register, which means a macro can be constructed directly with \texttt{:let @q = '...'} rather than performed live and captured. This is useful whenever the desired keystroke sequence is easier to reason about or generate programmatically than to physically type and record: a macro assembled from a computed string, or one that needs to embed a literal escape or control character that would be awkward to record by hand (using \texttt{\textbackslash{}<Esc>} or similar notation within the assigned string to represent it), is often more reliably constructed this way than through live recording. This is also the natural bridge into Chapter~\ref{ch:vimscript}: a register assignment is itself a single line of Vimscript, and a macro built this way is, formally, no different from any other piece of scripted automation covered next.

\section{Editing a Recorded Macro}

Because a macro's contents are register text, they can be inspected and modified with the same tools that inspect and modify any other register content. \texttt{:let @q} (with no assignment) displays register \texttt{q}'s current contents in the command-line area, letting a recorded macro be reviewed after the fact. A more thorough editing workflow pastes the register's contents into the buffer as ordinary text (\texttt{"qp}, using the put command from Chapter~\ref{ch:registers}), edits that text with any of the commands covered throughout Parts II and III, and then yanks the corrected line back into the register with \texttt{"qy\$} rather than re-recording the entire macro from scratch. This is generally faster than live re-recording for anything beyond a trivial correction, and it is worth knowing as the default approach for fixing a macro that almost worked rather than one that failed outright.

\section{The Pain-Point Threshold}

The testing-before-scaling discipline described above is really an informal version of a more general question: given that recording, validating, and running a macro all cost time up front, at what point does that up-front cost pay for itself against simply performing the transformation by hand each time it is needed? Put concretely: if a manual repetition costs some amount of time and attention per instance, and recording a macro costs a roughly fixed amount of setup time regardless of how many instances it is eventually run against, then a macro is worth recording once the number of remaining instances is large enough that the fixed setup cost, spread across all of them, comes out lower than doing each one by hand. For three or four remaining instances, that threshold is very often not met, and manual editing remains the faster choice; for fifty structurally consistent instances, it very clearly is. Reasoning explicitly in these terms, rather than reaching for a macro reflexively the moment any repetition appears, is worth doing consciously until it becomes second nature, since a macro recorded (and then debugged, when it inevitably breaks on some unanticipated variation) against only three or four remaining instances frequently ends up costing more total time than it saved.

This same reasoning has a second edge worth stating explicitly: the threshold above assumes the repeated material is genuinely structurally consistent. A macro recorded against material that only looks consistent, but in fact varies in some way the macro's author did not anticipate, will fail partway through a large replay, and the cost of noticing the failure, diagnosing what varied, and cleaning up whatever the macro did to the instances it mishandled can easily exceed whatever time the macro saved on the instances it handled correctly. This is the deeper reason behind the testing discipline from earlier in this chapter: validating against a few instances before scaling up is not merely a way to catch an outright broken macro, it is a way of empirically checking whether the structural consistency the whole cost-benefit calculation depends on actually holds, before committing to it at scale.

\section{Macro Libraries}

A macro recorded during one editing session does not persist to the next unless deliberately saved, since registers, like the undo history covered in Chapter~\ref{ch:undo}, are session-local by default. Where a particular transformation is needed repeatedly across many separate sessions, rather than repeatedly within a single one, the more durable solution is to store its \texttt{:let @q = '...'} definition in the vimrc introduced in Chapter~\ref{ch:installing}, populating register \texttt{q} (or whichever register is chosen) automatically every time Vim starts. A reader who finds themselves recording substantially the same macro in separate sessions, days or weeks apart, has outgrown live recording for that particular transformation and should consider promoting it into the vimrc as a standing, always-available macro instead, which is a small-scale instance of the same promotion-from-manual-to-persistent-automation pattern that motivates moving from a macro to a full Vimscript function, the subject of the next chapter, once a transformation's complexity outgrows what a macro can comfortably express.
